As dealing with vector fields on manifolds, we will work exclusively with Lie algebras over reals.
Probably the easiest way of constructing a nilpotent Lie algebra of vector fields is to consider vector fields with negative degree with respect to a positive dilation on ${\mathbb R}^n$. To be more precise, let us associate with the coordinates $\big(x^1,\dots,x^n\big)$ positive integer degrees (weights) $(w_1,\dots,w_n)$. This is equivalent to picking up
the \emph{weight vector field}
\[
\nabla^h=\sum_iw_ix^i\partial_{x^i} ,
\] or a one-parametric family of positive \emph{dilations} (called in \cite{Bruce:16, Grabowski:2006,Grabowski:2012}, a \emph{homogeneity structure})
\[
h_t\big(x^1,\dots,x^n\big)=\big(t^{w_1}x^1,\dots,t^{w_n}x^n\big) .
\]
The biggest $w_i$ is called the \emph{degree} of the homogeneity structure. The weight vector field defines \emph{homogeneous functions of degree~$k$} as those $f\in C^\infty\big({\mathbb R}^n\big)$
for which $\nabla^h(f)=kf$. One can prove~\cite{Grabowski:2012} that only non-negative homogeneity degrees are allowed and that each homogeneous function is polynomial in $\big(x^1,\dots,x^n\big)$. This can be extended to tensors, e.g., a~vector field $X$ is homogeneous of degree $k$ (what we denote by $X\in\mathfrak{g}^k(h)$) if $\big[\nabla^h,X\big]=kX$. This time, $k$ may be negative, for instance $\deg\big(\partial_{x^i}\big)=-w_i$ while $\deg\big(x^j\partial_{x^i}\big)=w_j-w_i$.
Since $\deg(X)\ge -w(h)=\min\{-w_1,\dots,-w_n\}$ and, for homogenous $X$, $Y$,
\[ \deg([X,Y])=\deg(X)+\deg(Y) ,\]
the family of vector fields of negative degrees
\[ \mathfrak{g}^{<0}(h)=\bigoplus_{i=-w(h)}^{-1}\mathfrak{g}^i(h)\]
is a finite-dimensional nilpotent Lie algebra (cf.~\cite{JG90,MK88}).
This (except for finite-dimensionality) remains valid for general graded supermanifolds~\cite{Voronov10}, in particular graded bundles in the sense of~\cite{Grabowski:2012}.

It is an interesting observation that any transitive nilpotent Lie algebra $L$ of vector fields defined on a neighbourhood of $0\in{\mathbb R}^n$ is locally a subalgebra of $\mathfrak{g}^{<0}(h)$ for some dilation $h$ on~${\mathbb R}^n$ \cite{JG90} (cf.\ also~\cite{MK88}, where the assumptions are a little bit stronger). The nilpotent algebra~$L$ is therefore automatically finite-dimensional and polynomial in appropriate homogeneous coordinates $\big(x^1,\dots,x^n\big)$. The homogeneity structure is uniquely defined by the structure of $L$, so that $\mathfrak{g}^{<0}(h)$ can be viewed as nilpotent `enveloping' or `prolongation' of $L$. Note also that every finite-dimensional (real) Lie algebra is a transitive Lie algebra of vector fields as the Lie algebra of left-invariant vector fields on the corresponding 1-connected Lie group.

The homogeneity structures defined in \cite{Grabowski:2012} are a little bit more general: they admit coordinates of degree $0$, so that the manifold is a fibration over the quotient manifold corresponding to zero-degree coordinates (graded bundle). In particular, graded bundles of degree 1 are just vector bundles.
The family of vector fields $\mathfrak{g}^{<0}(h)$ is still nilpotent, but it is a subalgebra in the family of vector fields with non-positive degrees,
\[ \mathfrak{g}^{\le 0}(h)=\bigoplus_{i=-w(h)}^{0}\mathfrak{g}^i(h) .\]
The latter is generally not solvable, but any finite-dimensional Lie algebra $L$ in $\mathfrak{g}^{\le 0}(h)$ such that $[L,L]\subset\mathfrak{g}^{< 0}(h)$ is solvable.
This is because, for finite-dimensional Lie algebras, $L$ is solvable if and only if $[L,L]$ is nilpotent (this is not true in infinite dimensions). In other words the \emph{nilradical} $\mathfrak{nr}(L)$ of $L$ contains $[L,L]$. Solvable Lie algebras of vector fields in $\mathfrak{g}^{\le 0}(h)$ we will call \emph{dilational} (with respect to~$h$).

If we distinguish coordinates $x^1,\dots,x^p$ of zero degree (weight) and coordinates of positive degree (weight) $z^1_{w_1},\dots,z^q_{w_q}$, where $w_i>0$ indicates the weight, then locally the homogeneity structure (dilation)
looks like (see \cite{Grabowski:2012})
\[ h_t(x,z)=\big(x^1,\dots,x^p,t^{w_1}z^1_{w_1},\dots,t^{w_q}z^q_{w_q}\big)\]
and smooth functions of degree $\le w$ are polynomials with coefficients in functions depending on~$x$:
\[ f(x,y)=\sum_{|\alpha|\le w}f_\alpha(x)\big(z_{w_1}^1\big)^{\alpha_1}\cdots \big(z_{w_q}^q\big)^{\alpha_q},
\]
where for $\alpha\in{\mathbb N}^q$, we put $|\alpha|=\sum_iw_i\alpha_i$. The method of defining dilational transitive solvable Lie algebras of vector fields is illustrated by the following example.

\begin{Example}\label{1.1}
On ${\mathbb R}^4$ with local coordinates $(x,y,z,u)$ we define $h$ declaring that $x$ is of weight~0, $y$ is of weight 1, $z$ is of weight~2, $u$ is of weight~3. The nilpotent part $\mathfrak{g}^{< 0}(h)$ is a module over ${\mathbb A}=C^\infty({\mathbb R})=\{ f(x)\}$ generated by polynomial vector fields
\begin{gather*} \mathfrak{g}^{-3}(h)={\mathbb A}[ \partial_u ], \qquad \mathfrak{g}^{-2}(h)={\mathbb A}[ y\partial_u, \partial_z ], \qquad
\mathfrak{g}^{-1}(h)={\mathbb A}\big[ z\partial_u, y^2\partial_u, y\partial_z, \partial_y \big].
\end{gather*}
We can now choose a commutative algebra of vector fields of weight 0 completing the above to a transitive algebra,
e.g.,
\[ \partial_x, \ y\partial_y, \ z\partial_z, \ u\partial_u ,\]
 and choose dilational solvable transitive Lie algebras of vector fields as transitive
subalgebras of
\begin{gather}\label{ida}\langle \partial_x, y\partial_y, z\partial_z, u\partial_u\rangle\oplus{\mathbb A}\left[ \partial_u, y\partial_u, \partial_z, \partial_u, z\partial_u, y^2\partial_u, y\partial_z, \partial_y \right].
\end{gather}
\end{Example}

Our aim in this note is to show that any transitive solvable finite-dimensional Lie algebra~$L$ of vector fields, defined locally in a~neighbourhood of $0\in{\mathbb R}^n$, is dilational, i.e., it is a~subalgebra of~$\mathfrak{g}^{\le 0}(h)$, with $[L,L]\subset\mathfrak{g}^{< 0}(h)$, for a homogeneity structure $h$ on ${\mathbb R}^n$ and appropriate homogeneous coordinates.
This is of course far from fully classifying solvable Lie algebras of vector fields, but it gives a nice structural information about $L$. This also allows for an easy way of constructing solvable Lie algebras of vector fields. Note also that our methods are constructive and serve as well for a large variety of infinite-dimensional Lie algebras of vector fields, especially in the analytical context.

We prove also the Lie's conjecture for solvable Lie algebras. The conjecture states that for any finite-dimensional complex transitive Lie algebra $L$ of vector fields in ${\mathbb R}^n$ one can find coordinates $y^1,\dots ,y^n$ in which all coefficients of vector fields from $L$ lie in the algebra generated by the~$y^i$ and the exponentials $\exp\big(\lambda y^i\big)$.

\section{Nilpotent Lie algebras of vector fields}
Our main result slightly depends on the fact whether $L$ itself is nilpotent or not. If yes, then we recover the result proved already in~\cite{JG90}. We will reformulate this result as follows.

The following is well known. Let $L$ be a Lie algebra. We define the lower central series of $L$ inductively:
\[ L=L^1 ,\qquad L^{i+1} = \big[L,L^i\big] \qquad \text{for}\quad i\ge 1 .\]
\begin{Definition} If the lower central series \emph{stops at zero}, i.e.,
$L^{k+1}=\{ 0\}$ for some $k$, then we say that~$L$ is \emph{nilpotent}.
The smallest such $k$ we call the \emph{height} of the nilpotent Lie algebra.
\end{Definition}
Let us recall from the introduction that on ${\mathbb R}^n$ we consider a positive dilation, i.e., an action~$h$ of the multiplicative monoid $({\mathbb R},\cdot)$
of the form
\[ h_t\big(x^1,\dots,x^n\big)=\big(t^{w_1}x^1,\dots,t^{w_n}x^n\big) ,\qquad t\in{\mathbb R} ,\]
where $w_i\in{\mathbb Z}$, $w_i>0$, $w(h)=\max\{w_1,\ldots,w_n\}.$ The associated \emph{weight vector field} reads $\nabla^h=\sum_iw_ix^i\partial_{x^i}$.
A smooth function $f$ defined in a neighbourhood of $0\in{\mathbb R}^n$ is homogeneous of degree $a$ if
$\nabla^h(f)=af$. Note that only non-negative integer homogeneity degrees are possible.
Similarly, a vector field $X$ is homogeneous of degree $a$ if
$\big[\nabla^h,X\big]=aX$, however now negative degrees not less than $-w(h)$ are allowed.
The family of homogeneous vector fields (of degree~$a$) is denoted by~$\mathfrak{g}^a(h)$.

We observe that the family of vector fields of negative degrees
\[ \mathfrak{g}^{<0}(h)=\bigoplus_{i=-w(h)}^{-1}\mathfrak{g}^i(h)\]
is a transitive (finite-dimensional) nilpotent Lie algebra of vector fields (cf.~\cite{JG90,MK88}).
The transitivity means that the vector fields of negative degrees span ${\mathsf T} {\mathbb R}^n$.
Of course, any transitive subalgebra of $\mathfrak{g}^{<0}(h)$ is again transitive nilpotent Lie algebra of vector fields. We can, of course, restrict $\mathfrak{g}^{<0}(h)$ to a neighbourhood $U$ of $0\in{\mathbb R}^n$ which gives a Lie algebra $\mathfrak{g}^{<0}(U,h)$ isomorphic to $\mathfrak{g}^{<0}(h)$. In this case we will speak about a \emph{local Lie algebra of vector fields}.
\begin{Example}
Consider the dilation $h$ on ${\mathbb R}^2=\{(y,z)\}$ such that $y$ is of degree 1 and $z$ is of degree 2. The Lie algebra of vector fields with negative degrees is spanned by $\langle\partial_y,\partial_z,y\partial_z\rangle$ which is transitive nilpotent of height 2 since $L^2=[L,L]=\langle \partial_z\rangle$.
\end{Example}

For a nilpotent local Lie algebra $L$ of vector fields on ${\mathbb R}^n$ of height $k$ we consider the descending series of finite-dimensional vector spaces $L^i(0)\subset{\mathsf T}_0{\mathbb R}^n$, spanned by values at $0\in{\mathbb R}^n$ of vector fields belonging to $L^i$,
\begin{gather*}%\label{tan}
L(0)=L^1(0)={\mathsf T}_0{\mathbb R}^n\supset L^2(0)\supset \cdots\supset L^{k}(0)\supset L^{k+1}(0)=\{ 0\} .
\end{gather*}
Let us rewrite this sequence, collecting terms of the same dimension of the space and indicating this dimension by the lower indices:
\begin{gather*}
 \underbrace{L^1_n(0)=\cdots= L^{r_1}_n(0)}_{r_1}\supsetneq
\underbrace{L^{r_1+1}_{n-b_1}(0)=\cdots= L^{r_2}_{n-b_1}(0)}_{r_2-r_1}\supsetneq
\cdots\supsetneq\\
\qquad{} \supsetneq\underbrace{L^{r_{m-1}+1}_{n-b_{m-1}}(0)=\cdots= L^{r_m}_{n-b_{m-1}}(0)}_{r_m-r_{m-1}}
\supsetneq \{ 0\} .
\end{gather*}
Put $r_0=0$, $b_0=0$, $a_0=0$, $b_m=n$, and $a_i=b_i-b_{i-1}$ for $i>0$. Of course, $a_1+\cdots+a_m=n$.
The main result in \cite{JG90} can be reformulated as follows.
\begin{Theorem}\label{nil}
Let $L$ be a transitive local nilpotent Lie algebra of vector fields on ${\mathbb R}^n$ and let~$r_i$,~$b_i$ and~$a_i$, $i=0,\dots,m$, be defined as above. There exist local variables $\big(y^1,\dots,y^n\big)$ in a~neighbourhood $U$ of $0\in{\mathbb R}^n$ such that $y^1,\dots, y^{b_j}$ vanish on $L^{r_j+1}$ and $\partial_{y^{b_j+1}},\dots,\partial_{y^{b_{j+1}}}$ span $L^{r_{j}+1}(0)/L^{r_{j+1}}(0)$, $j=1,\dots,m-1$. Let $h$ be a positive dilation of ${\mathbb R}^n$ such that the first $a_1$ variables, $y^1,\dots,y^{a_1}$, are of degree $r_1>0$, the next $a_2$ variables are of degree~$r_2$, etc., up to final~$a_m$ variables, $y^{b_{m-1}+1},\dots,y^n$ which are of degree $r_{m}$.
Then, with respect to this dilation, $L\subset\mathfrak{g}^{<0}(U,h)$ and
\[ L^j\subset \bigoplus_{i=-k}^{-j}\mathfrak{g}^i(U,h) ,\]
where $k$ denotes the height of~$L$.
In particular, the algebra $L$ is finite-dimensional and polynomial in the chosen coordinates.
\end{Theorem}
\begin{Remark}
Note that the dilation above, so $\mathfrak{g}^{<0}(U,h)$, is completely determined (up to a~diffeomorphism) by the dimensions of $L^i(0)$, so the algebra of vector fields $L$.
\end{Remark}
\begin{Example}
Consider the Lie algebra of vector fields on ${\mathbb R}^2$ spanned by
$L=L^1=\big\langle\partial_x,\partial_y,y\partial_x,\allowbreak y^2\partial_x\big\rangle$. It is nilpotent of height 3: $L^2=\langle\partial_x,y\partial_x\rangle$, $L^3=\langle\partial_x\rangle$. According to the theorem, we associate with $y$ degree 1 and with $x$ degree 3, so $L$ is spanned by vector fields of degree $\le -1$.
\end{Example}
\begin{Remark} Let us now consider nilpotent Lie algebra $L$ included in
\[ \mathfrak{g}^{<0}(h_0)=\bigoplus_{i=-w(h_0)}^{-1}\mathfrak{g}^i(h_0)\]
for a given dilation $h_0$. We observe that the dilation $h$ determined by $L$ according to Theorem~\ref{nil} need not to recover $h_0$, so in general $h\ne h_0$. For instance if coordinate $x$ is of degree 0 and $y$ is of degree $2$ for $h_0$, for the commutative algebra $L=\langle\partial_x,\partial_y\rangle$ it is sufficient to take $h$ for which~$x$ and~$y$ are of degree~1.
\end{Remark}

\section{Solvable algebras: descending series of ideals and foliations}
Let $L$ be a Lie algebra. We define the \emph{nilradical series} of $L$ inductively:
\begin{gather}\label{hcs} L^0=L ,\qquad L^1=\mathfrak{nr}(L) , \qquad L^{i+1}=\big[\mathfrak{nr}(L),L^i\big]\qquad \text{for}\quad i\ge 1 ,
\end{gather}
where $\mathfrak{nr}(L)$ is the nilradical of $L$.
Its name will be justified by the fact that it its main part (i.e., starting from $i=1$) is the lower central series of the nilradical of $L$. The following observation is fundamental. Note a shift in the notation in comparison with the nilpotent case due to the fact that we will have to accept coordinates of degree~$0$.
\begin{Proposition}
For a finite-dimensional $L$ $($over a field of characteristic~$0)$, the nilradical series is central, i.e., $[L,L]\subset\mathfrak{nr}(L)$, if and only if~$L$ is solvable.
\end{Proposition}
\begin{proof}
It is easy to see that the nilradical series is a central series if and only if $[L,L]$ is nilpotent. For finite-dimensional Lie algebras (over a field of characteristic~$0$) this is equivalent to solvability.
\end{proof}
\begin{Remark}
As the main result we are approaching is known to be valid for nilpotent Lie algebras of vector fields, we can in general assume that $L^0=L\ne L^1=\mathfrak{nr}(L)$. It is also obvious that the series (\ref{hcs}) stops at zero, i.e., $L^{k+1}=\{ 0\}$ for some non-negative integer~$k$.
\end{Remark}
\begin{Example} The proposition is not true in infinite-dimension. For instance, consider the Lie algebra $L$ spanned by elements $\langle x,a,h_1,h_2,\dots\rangle$ subject to the Lie bracket
\[ [x,a]=a ,\qquad [x,h_i]=ih_i ,\qquad [a,h_i]=h_{i+1} ,\qquad [h_i,h_j]=0 .\]
The derived ideal $[L,L]$ is spanned by $\langle a,h_1,h_2,\dots\rangle$ and it is solvable but not nilpotent.
\end{Example}

Now, we will realize a (real) finite-dimensional solvable Lie algebra $L$ as a \emph{local transitive Lie algebra of vector fields on ${\mathbb R}^n$}, i.e., a Lie algebra of smooth vector fields defined in a neighbourhood $U$ of $0\in{\mathbb R}^n$ such that the vector fields from $L$ span the whole tangent space ${\mathsf T}_0{\mathbb R}^n$, i.e.,
\[ L(0)={\mathsf T}_0{\mathbb R}^n\simeq{\mathbb R}^n .\]
We can of course equivalently assume that $L$ spans the tangent spaces ${\mathsf T}_x{\mathbb R}^n$ for all $x\in U$.
As the generalized distributions $V^i$ generated in $U$ by $L^i$ are invariant with respect to the infinitesimal action of~$L$, they are invariant with respect to the local Lie group corresponding to $L$ which transitively acts by local diffeomorphisms in a neighbourhood of $0\in{\mathbb R}^n$, say $U$ itself (cf.~\cite{Palais}). The distributions are therefore regular and clearly involutive, so define foliations ${\mathcal F}^i$ on $U$. Thus we get the following.

\begin{Proposition}
If $L$ is a local transitive finite-dimensional solvable Lie algebra of vector fields in a neighbourhood of $0\in{\mathbb R}^n$, then ideals $L^i$ from the nilradical series
\[ L=L^0\vartriangleright L^1\vartriangleright L^2\vartriangleright\cdots\vartriangleright L^k\vartriangleright L^{k+1}=\{ 0\}\]
of $L$ generate a descending series
\begin{gather}\label{fol}{\mathcal F}^0=\{ U\}\supset{\mathcal F}^1\supset{\mathcal F}^2\supset \cdots\supset {\mathcal F}^{k}\supset {\mathcal F}^{k+1}=U \end{gather}
of foliations in a neighbourhood $U$ of $0\in{\mathbb R}^n$.
\end{Proposition}

Inclusion ${\mathcal F}^i\supset{\mathcal F}^{i+1}$ means that leaves of foliation ${\mathcal F}^{i+1}$ are submanifolds of the leaves of foliation ${\mathcal F}^{i}$. Note that the series of foliations~(\ref{fol}) need not be strictly decreasing and that we understand that $L^k\ne\{ 0\}$.
\begin{Example}\label{ex2}
Let $L$ be the solvable Lie algebra of vector fields on ${\mathbb R}^2$ spanned by
\[ L=\big\langle \partial_x ,\partial_y ,x\partial_x ,y\partial_y ,y^2\partial_x ,y\partial_x\big\rangle .\]
Then,
\[ L^1=[L,L]=\big\langle \partial_x ,\partial_y ,y^2\partial_x ,y\partial_x\big\rangle ,\]
so ${\mathcal F}^1={\mathcal F}^0=\big\{{\mathbb R}^2\big\}$. Moreover, $L^2=\big[L^1,L^1\big]=\langle\partial_x ,y\partial_x\rangle$ and $L^3=\big[L^1,L^2\big]=\langle\partial_x\rangle$, so ${\mathcal F}^2={\mathcal F}^3$ are generated by $\partial_x$ and $k=3$.
\end{Example}
\begin{Remark}We consider only transitive Lie algebras, since non-transitive ones are much harder to describe. In particular, in the non-transitive case we can have huge infinite-dimensional commutative Lie algebras of vector fields, as for example
the algebra $L=\langle f(y)\partial_x\rangle$ on ${\mathbb R}^2$. Here, $f$ is an arbitrary smooth function of one variable.
Any full description of such algebras seems extremely difficult.
\end{Remark}

Let us fix a local solvable finite-dimensional algebra of vector fields and pick up from the sequence (\ref{fol}) the subsequence
\begin{gather}\label{fol1}
{\mathcal F}^{r_1}\supsetneq{\mathcal F}^{r_2}\supsetneq \cdots\supsetneq {\mathcal F}^{r_m}\supsetneq {\mathcal F}^{r_{m+1}}=U \end{gather}
consisting of those ${\mathcal F}^i$ for which the dimension of ${\mathcal F}^{i+1}$ is smaller than that of ${\mathcal F}^i$. In particular, ${\mathcal F}^{r_m}={\mathcal F}^k$. Put
\begin{gather}\label{adef}a_i=\dim\big({\mathcal F}^{r_i}\big)-\dim\big({\mathcal F}^{r_i+1}\big) ,\qquad i=1,\dots,m .\end{gather}
Of course, $a_i>0$, $a_m=\dim\big({\mathcal F}^k\big)$, and $a_1+\cdots+a_m=n$, so that $0=b_0<b_1<b_2<\cdots <b_m=n$, where $b_j=a_1+\cdots+a_j=\dim\big({\mathcal F}^{r_j}\big)$ for $j>0$.
Hence, the series~(\ref{fol}) looks like
\begin{gather*}%\label{fol3}
\underbrace{{\mathcal F}^{0}_n=\cdots={\mathcal F}^{r_1}_n}_{r_1+1}\supsetneq
\underbrace{{\mathcal F}^{r_1+1}_{n-b_1}=\cdots={\mathcal F}^{r_2}_{n-b_1}}_{r_2-r_1}\supsetneq
\cdots\supsetneq
\underbrace{{\mathcal F}^{r{_m-1}+1}_{n-b_{m-1}}=\cdots={\mathcal F}^{r_m}_{n-b_{m-1}}}_{r_m-r_{m-1}}
\supsetneq U ,
\end{gather*}
where with the lower indices we indicated the dimension of the leaves of the foliation. In the above example it is a subsequence ${\mathcal F}^1$, ${\mathcal F}^3$, so that $a_1=a_2=1$.

\begin{Proposition}[cf.~\cite{CFG16, CFGR15}]\label{p1}
There are coordinates $y^1,\dots,y^n$ in a neighbourhood of $0\in{\mathbb R}^n$ such that:
\begin{enumerate}\itemsep=0pt
\item[$1.$] The coordinates $y^1,\dots, y^{b_j}$ define ${\mathcal F}^{r_{j+1}}$ as their level sets for each $j=1,\dots,m$.
\item[$2.$] There are $Y_1,\dots,Y_n\in L$ such that, for each $j=1,\dots,m$, the vector fields $Y_{b_{j-1}+1},\dots,Y_n$ generate ${\mathcal F}^{r_j}$ and
\begin{gather}\label{gen} Y_k=\partial_{y^k} \qquad \operatorname{mod}{\mathcal F}^{r_{j+1}} ,\end{gather}
for $k=b_{j-1}+1,\dots,b_j$.
\item[$3.$] The generators of ${\mathcal F}^k$ are the vector fields
\begin{gather}\label{fine} Y_{b_{m-1}+1}=\partial_{y^{b_{m-1}+1}} ,\qquad \dots,\qquad Y_n=\partial_{y^n}\end{gather}
belonging to $L^k$.
\item[$4.$] Moreover, every $X\in L$ can be written as
\begin{gather}\label{form}X=\sum_{j=b_{m-1}+1}^n\left(\sum_{i=b_{m-1}+1}^n f_i^j\big(y^1,\dots,y^{b_{m-1}}\big)y^i+g^j\big(y^1,\dots,y^{b_{m-1}}\big)\right)\partial_{y^j}
+\bar{X} ,
\end{gather}
where the vector field $\bar{X}$ depends on the first $b_{m-1}$ variables only,
\begin{gather*}%\label{form1}
\bar{X}=\sum_{i=1}^{b_{m-1}}v^i\big(y^1,\dots,y^{b_{m-1}}\big)\partial_{y^i} .\end{gather*}
If $X\in L^1$, then
\begin{gather*}%\label{form2}
X=\sum_{j=b_{m-1}+1}^nu^j\big(y^1,\dots,y^{b_{m-1}}\big)\partial_{y^j}
+\bar{X} .
\end{gather*}

\item[$5.$] The vector fields $\bar{X}$ form a local Lie algebra $\bar{L}$ of vector fields in a neighbourhood $\bar{U}$
of $0\in{\mathbb R}^{b_{m-1}}$ and $L\ni X\mapsto \bar{X}\in\overline{L}$ is a Lie homomorphism onto $\bar{L}$. The algebra $\bar{L}$ is therefore solvable and transitive finite-dimensional Lie algebra of vector fields in $b_{m-1}=k-a_m$ variables.
\end{enumerate}
\end{Proposition}
\begin{proof} The coordinates $y^1,\ldots, y^n$ will be constructed inductively. In the following $V^i$ denotes the distribution spanned by elements of $L^i$ and $L^i(0)=V^i\cap{\mathsf T}_0 U$. In particular $V^{r_1}={\mathsf T} U$ since~$L$ is transitive.
Take $Y_1,\dots,Y_{b_1}\in L^{r_1}$ such that $Y_1(0),\dots,Y_{b_1}(0)$ generate $L^{r_1}(0)$ $\operatorname{mod} L^{r_2}(0)$. Hence, $Y_1,\dots,Y_{b_1}$ span $V^{r_1}$ $\operatorname{mod} V^{r_2}$ in a neighbourhood of $0\in{\mathbb R}^n$.
Let now $\alpha^1,\dots,\alpha^{b_1}$ be 1-forms in a neighbourhood on $0\in{\mathbb R}^n$, vanishing on $V^{r_2}$ and such that
$\alpha^l(Y_s)=\delta^l_s$, $l,s=1,\dots, b_1$. These 1-forms are closed. Indeed,
\[ {\rm d}\alpha^l(Y_s,Y_{s'})=Y_s\big(\alpha^l(Y_{s'})\big)- Y_{s'}\big(\alpha^l(Y_{s})\big)-\alpha^l([Y_s,Y_{s'}])=0 .\] Note that $\alpha^l([Y_s,Y_{s'}])=0$, since
\[ [Y_s,Y_{s'}]\subset \big[L^{r_1},L^{r_1}\big]\subset L^{r_1+1},\]
so that $[Y_s,Y_{s'}]$ is a section of the distribution $V^{r_1+1}=V^{r_2}$.
Similarly, for $X,X'\in L^{r_2}$ we have
\[ {\rm d}\alpha^l(Y_s,X)=0 \qquad\text{and}\qquad {\rm d}\alpha^l(X,X')=0 .\]
As $L^{r_2}$ spans $V^{r_2}$, this completes the proof that ${\rm d}\alpha^l=0$.
We can therefore (uniquely) choose functions $y^1,\dots,y^{b_1}$ vanishing at $0\in{\mathbb R}^n$ such that $\alpha^l={\rm d} y^l$. In particular, $y^j$ are constant on leaves of ${\mathcal F}^{r_2}$, $Y_s\big(y^l\big)=\delta^l_s$. Note that $u=\big(y^1,\ldots, y^{b_1}\big)$ parameterize leaves of ${\mathcal F}^{r_2}$. By $F(u)$ we will denote the leaf of ${\mathcal F}^{r_2}$ corresponding to $u$.

Similarly, we can choose $Y_{b_1+1},\dots,Y_{b_2}\in L^{r_2}$ which span $V^{r_2} \operatorname{mod} V^{r_3}$ in a neighbourhood of $0\in{\mathbb R}^n$. As before, on each leaf $F(u)$ of ${\mathcal F}^{r_2}$ we can find coordinates $y^{b_1+1}_u,\dots,y_u^{b_2}$ which are constant on the leaves of ${\mathcal F}^{r_3}_{|F(u)}$ and such that, on $F(u)$, we have $Y_s\big(y^l_u\big)=\delta^l_s$. The functions $y^{b_1+1}_u,\dots,y_u^{b_2}$ are determined up to a constant on each leaf. If we choose them so that they vanish on a smooth curve through $0\in{\mathbb R}^n$ and transversal to the leaves of ${\mathcal F}^{r_2}$, they smoothly depend on $u$ and define smooth functions $y^{b_1+1},\dots,y^{b_2}$ in a neighbourhood of $0\in{\mathbb R}^n$.

At this point of the construction we have vector fields $Y_1,\dots,Y_{b_2}$ and smooth functions $y^1,\dots,y^{b_2}$. The functions $\big(y^1,\dots,y^{b_1}\big)$ define the foliation ${\mathcal F}^{r_2}$, and ($y^1,\dots,y^{b_2}$) define ${\mathcal F}^{r_3}$. Moreover, $Y_s\big(y^l\big)=\delta^l_s$ for $s,l=1,\dots,b_1$ and for $s,l=b_1+1,\dots,b_2$.
Proceeding inductively in this way, we finally get vector fields $Y_1,\dots,Y_n\in L$ and coordinates $y^1,\dots,y^n$ in a neighbourhood of $0\in{\mathbb R}^n$ with the desired properties.

Finally, from (\ref{gen}) we deduce (\ref{fine}), i.e., $V^{k}=V^{r_m}$ is generated by $Y_{b_{m-1}+1}=\partial_{y^{b_{m-1}+1}}$ up to $Y_n=\partial_{y^n}$. Hence, every $Y\in L^k$ is of the form
\begin{gather}\label{last} Y=\sum_{j=b_{m-1}+1}^n\gamma_j(y)\partial_{y^j} .\end{gather}
But $L^k$ is commutative, so that $[\partial_{y^i},Y]=0$ for $i=b_{m-1}+1,\dots,n$, that implies that $\gamma_j$ depends on $\big(y^1,\ldots,y^{b_{m-1}}\big)$ only. Since, $L^k$ is an ideal, for every $X\in L$
we have that $[\partial_{y^i},X]$ is of the form~(\ref{last}). This, in turn, implies~(\ref{form}) and \begin{gather*}%\label{mor}
\overline{[X,Y]}=[\overline{X},\overline{Y}] .\tag*{\qed}\end{gather*}\renewcommand{\qed}{}
\end{proof}

\begin{Example}\label{ex3}
Let $n=1$, which means we are looking for local finite-dimensional solvable Lie algebras $L$ of vector fields on the real line.
As ${\mathcal F}^k=\{{\mathbb R}\}$, the algebra $L^k$ contains $Y_1=\partial_y$. As there are no other variables, according to (\ref{form}),
each $X\in L$ is locally of the form $X=(ay+b)\partial_y$. Hence, $L=\langle\partial_y\rangle$ or $L=\langle\partial_y ,y\partial_y\rangle$. Thus we recovered the classical result known already to S.~Lie. In the first case of $L=\langle\partial_y\rangle$ the algebra is nilpotent, so in view of Theorem~\ref{nil} we associate the degree~$1$ to the variable~$y$. In the second case of $L=\langle\partial_y\rangle$ which is solvable but not nilpotent we also associate the degree~$1$ to the variable~$y$ in preparation for the next section. We have then~$\partial_y$ of degree~$-1$ and $y\partial_y$ of degree~$0$.
\end{Example}

\section{Solvable Lie algebras are dilational}
\begin{Theorem}\label{main}
Let $L$ be a finite-dimensional transitive local solvable Lie algebra of vector fields on ${\mathbb R}^n$ and let $r_i$ and $a_i$, $i=0,\dots,m$, be defined as in \eqref{fol1} and \eqref{adef}. Then, there is a non-negative dilation $($homogeneity structure$)$ on ${\mathbb R}^n$ which can be described in a constructive way such that the first $a_1$ variables are of degree $r_1$, the next $a_2$ variables are of degree $r_2$, etc., up to final $a_m$ variables of degree $r_m$, completely determined $($up to a diffeomorphism$)$ by~$L$, so the eigenvalues and multiplicities of the dilation have invariant meaning.
Moreover, $L\subset\mathfrak{g}^{\le 0}(U,h)$ and
$$L^j\subset \bigoplus_{i=-k}^{-j}\mathfrak{g}^i(U,h) .$$
Actually, there is a finite dimensional commutative subalgebra $\mathfrak{{g}}^{0}_0(U,h)$ in $\mathfrak{{g}}^{0}(U,h)$ such that
\begin{gather}\label{solv} L \subset \mathfrak{g}^{<0}(U,h)\oplus \mathfrak{g}_0^{0}(U,h) .
\end{gather}
\end{Theorem}
\begin{proof} We shall proceed inductively with respect to $n$. The theorem is true for $n=1$ (see Example \ref{ex3}). Suppose it is true for $n-1$, $n\ge 2$. Let $\big(y^1,\dots,y^n\big)$ be variables described by Proposition \ref{p1}, let $h$ be a non-negative dilation (homogeneity structure) on ${\mathbb R}^n$ such that the first $a_1$ variables $y^1,\dots,y^{a_1}$ are of degree $r_1$, the next $a_2$ variables are of degree $r_2$, etc., up to final $a_m$ variables of degree $r_{m}$. Forgetting the last $a_m$ variables (i.e. passing to the manifold ${\mathbb R}^n/{\mathcal F}^k\subset {\mathbb R}^{b_{m-1}}$ of leaves of the foliation ${\mathcal F}^k$), we obtain a non-negative dilation $\bar{h}$ of ${\mathbb R}^{b_{m-1}}$ and Lie algebra $\bar{L}$ described in item~5 of Proposition~\ref{p1}. By the inductive assumption,
\begin{gather*}%\label{solv1}
\bar{L} \subset \mathfrak{g}^{<0}\big(\bar{U},\bar{h}\big)\oplus \mathfrak{g}_0^{0}\big(\bar{U},\bar{h}\big)
\end{gather*}
and
\[ \bar{L}^j\subset \bigoplus_{i=-k+1}^{-j}\mathfrak{g}^i\big(\bar{U},\bar{h}\big) ,\qquad j=1,\dots,k-1 .\]
Now, in view of Proposition \ref{p1}, every element of $L$ is of the form
\begin{gather}\label{form4} X=\sum_{j=b_{m-1}+1}^n\left(\sum_{i=b_{m-1}+1}^nf_i^j\big(y^1,\dots,y^{b_{m-1}}\big)y^i +g^j\big(y^1,\dots,y^{b_{m-1}}\big)\right)\partial_{y^j}
+\bar{X} ,
\end{gather}
while elements of $L^1$ are of the form
\begin{gather*} X=\sum_{j=b_{m-1}+1}^n u^j\big(y^1,\dots,y^{b_{m-1}}\big)\partial_{y^j}.
\end{gather*}
It suffices to show that functions $f_i^j\big(y^1,\dots,y^{b_{m-1}}\big)$ depend of coordinates of degree 0 only, so they are of degree zero, $g^j\big(y^1,\dots,y^{b_{m-1}}\big)$ and $u^j\big(y^1,\dots,y^{b_{m-1}}\big)$ are of degree $r_m$ and $r_m-1$ respectively.

Let us start with $L^1$. By definition it is a nilpotent Lie algebra. There are two possibilities: $r_1>0$, i.e., foliation ${\mathcal F}^1$ is trivial in a sense that it has only one leaf equal $U$, or $r_1=0$ so the foliation ${\mathcal F}^1$ is nontrivial. In the first case there are no variables of degree~$0$ and $L^1$ is a~transitive, nilpotent Lie algebra. Therefore we can use Theorem \ref{nil} according to which elements of $L^1$ are of degree at most $-1$, so functions $u^j\big(y^1,\dots,y^{b_{m-1}}\big)$ are of degree at most~\mbox{$r_m-1$}. In the second case $L^1$ is a transitive, nilpotent Lie algebra on each leaf of the foliation ${\mathcal F}^1$, i.e., $u_j\big(y^1,\dots,y^{b_{m-1}}\big)$ is of degree $r_m-1$ for each fixed $y^1=\alpha_1,\dots,y^{a_1}=\alpha_{a_1}$. But $y^1,\dots,y^{a_1}$ are of degree $0$ in this case, so $u_j\big(y^1,\dots,y^{b_{m-1}}\big)$ are of degree $r_m-1$ as a whole. Since variables $\big(y^{b_{m-1}+1},\ldots, y^n\big)$ are of degree $r_m$, this means that
\begin{gather*}%\label{ind1}
L^1\subset\mathfrak{g}^{<0}(U,h) .
\end{gather*}
Moreover, as $L$ is not nilpotent and $L^1\ne\{0\}$, elements of $L^k$ are of the form
\begin{gather}\label{form5} X=\sum_{j=b_{m-1}+1}^n u_j\big(y^1,\dots,y^{a_1}\big)\partial_{y^j} ,
\end{gather}
where $y^1,\dots,y^{a_1}$ are coordinates of degree 0 (or $u^j$ are constants when there are no coordinates of degree zero).

Take now $X\in L$ as given in (\ref{form4}). Since
\[ [\partial_{y^s},X]=\sum_{j=b_{m-1}+1}^nf_s^j\big(y^1,\dots,y^{b_{m-1}}\big)\partial_{y^j}\in L^k\]
for $s=b_{m-1}+1,\dots,n$, we get according to (\ref{form5}) that $f_s^j\big(y^1,\dots,y^{b_{m-1}}\big)$
are homogeneous of degree 0, i.e.,
\begin{gather*}%\label{L0}
f_s^j=f_s^j\big(y^1,\dots,y^{a_1}\big) \qquad \text{for} \quad s,j=b_{m-1}+1,\dots,n .\end{gather*}

What remains now is to prove that $g^j\big(y^1,\dots,y^{b_{m-1}}\big)$ are polynomial of degree $\le r_m$. We will do it inductively, proving that, for each $s=1,\dots,m$, the derivatives
$\partial_{y^r}g_j$ where $r=b_{s-1}+1,\dots,b_s$, are polynomial of degree $\le r_m-r_s$. This is trivially true for $s=m$,
as $g_j$ do not depend on variables $y^{b_{m-1}+1},\dots, y^n$. Suppose it is true for $s>1$, we will show that it true for $s-1$. Since $[\partial_{y^r},X]\in L^{r_{s-1}}$ for $r=b_{s-2}+1,\dots,b_{s-1}$ (recall that such $\partial_{y^r}\in L^{r_{s-1}} \operatorname{mod} L^{r_{s-1}+1})$, this implies that
\[\partial_{y^r}g_j\big(y^1,\dots,y^{b_{m-1}}\big)\]
are polynomials of degree $\le r_m-r_{s-1}$.
Now it is easy to see that, as $\partial_{y^r}g_j\big(y^1,\dots,y^{b_{m-1}}\big)$ is polynomial of degree $r_m-\deg(\partial_{y^r})$ for each $r$, the function $g_j\big(y^1,\dots,y^{b_{m-1}}\big)$ itself is polynomial of degree $\le r_m$.

Thus we have shown that $L\subset\mathfrak{g}^{\le 0}(U,h)$. Let $\mathfrak{g}^{0}_0(U,h)$ be the projection of $L$ onto $\mathfrak{g}^{0}(U,h)$. Then we have (\ref{solv}) and $\mathfrak{g}^{0}_0(U,h)$ is commutative.
Indeed, it follows from $[L,L]=L^1\subset\mathfrak{g}^{< 0}(U,h)$ and that $\mathfrak{g}^{0}(U,h)$ is a Lie algebra of vector fields. Since $L^1\subset \bigoplus_{i=-k}^{-1}\mathfrak{g}^i(U,h)$ and $L^j=\big[L^1,\big[L^1\big[\dots,\big[L^1,L^1\big]\dots\big]\big]\big]$ for $j>1$, we get
\[L^j\subset \bigoplus_{i=-k}^{-j}\mathfrak{g}^i(U,h) .\tag*{\qed}\]
\renewcommand{\qed}{}
\end{proof}
